\documentclass[11pt,a4paper]{article}
\usepackage[T1]{fontenc}
\usepackage[{{babel}}]{babel}
\usepackage{lmodern}
\usepackage[margin=2.3cm]{geometry}
\usepackage{amsmath, amssymb}
\usepackage[most]{tcolorbox}
\usepackage{enumitem}
\usepackage[hidelinks]{hyperref}

% {{fr:Encadrés du cours (numérotés par section)|en:Course boxes (numbered per section)}}
\newtcbtheorem[number within=section]{definition}{{{fr:Définition|en:Definition}}}%
  {colback=blue!4, colframe=blue!55!black, fonttitle=\bfseries}{def}
\newtcbtheorem[use counter from=definition]{theoreme}{{{fr:Théorème|en:Theorem}}}%
  {colback=red!4, colframe=red!55!black, fonttitle=\bfseries}{thm}
\newtcolorbox{exemple}{colback=green!4, colframe=green!45!black, title={{{fr:Exemple|en:Example}}}, fonttitle=\bfseries}
\newtcolorbox{attention}{colback=orange!6, colframe=orange!70!black, title={{{fr:Attention|en:Watch out}}}, fonttitle=\bfseries}

\title{{{title}}}
\author{{{author}}}
\date{{{institution}}}

\begin{document}
\maketitle
\tableofcontents

\section{{{fr:Les fonctions|en:Functions}}}

\begin{definition}{{{fr:Fonction|en:Function}}}{fonction}
  {{fr:Une fonction $f$ de $E$ dans $F$ associe à chaque $x \in E$ un unique $f(x) \in F$.|en:A function $f$ from $E$ to $F$ maps each $x \in E$ to a unique $f(x) \in F$.}}
\end{definition}

\begin{exemple}
  $f \colon x \mapsto x^2$ {{fr:est définie sur $\mathbb{R}$.|en:is defined on $\mathbb{R}$.}}
\end{exemple}

\begin{theoreme}{{{fr:Valeurs intermédiaires|en:Intermediate values}}}{tvi}
  {{fr:Si $f$ est continue sur $[a,b]$ et $f(a) \leq 0 \leq f(b)$, alors $f$ s'annule sur $[a,b]$.|en:If $f$ is continuous on $[a,b]$ and $f(a) \leq 0 \leq f(b)$, then $f$ vanishes on $[a,b]$.}}
\end{theoreme}

\begin{attention}
  {{fr:La continuité est indispensable (théorème~\ref{thm:tvi}).|en:Continuity is required (Theorem~\ref{thm:tvi}).}}
\end{attention}

\subsection*{Exercices}
\begin{enumerate}[label=\textbf{\arabic*.}]
  \item {{fr:Montrer que $x^3 + x - 1$ s'annule sur $[0,1]$.|en:Show that $x^3 + x - 1$ vanishes on $[0,1]$.}}
  \item {{fr:Donner un contre-exemple sans continuité.|en:Give a counter-example without continuity.}}
\end{enumerate}

\end{document}
